\documentclass[11pt]{article}
\usepackage[T1]{fontenc}
\usepackage[cmex10]{amsmath}
\usepackage{amssymb,amsthm}
\usepackage[hidelinks]{hyperref}
\usepackage[margin=1in]{geometry}
\DeclareFontShape{OMX}{cmex}{m}{n}{<->cmex10}{}

\newtheorem{theorem}{Theorem}
\newtheorem{lemma}{Lemma}
\DeclareMathOperator{\CT}{CT}
\DeclareMathOperator{\Arg}{arg}
\newcommand{\dd}{\mathrm d}

\title{Eventual real-rootedness of balanced array polynomials with fixed column count}
\hypersetup{pdftitle={Eventual real-rootedness of balanced array polynomials with fixed column count}}
\makeatletter
\renewcommand{\maketitle}{%
  \begingroup
  \centering\LARGE\@title\par
  \endgroup
  \vspace{0.4\baselineskip}}
\makeatother

\begin{document}
\maketitle

\begin{abstract}
We study the polynomial that counts binary arrays with equal column
sums and a prescribed bound on the row sums, weighted by the common
column sum. At any fixed column count and fixed proper row cap, we prove
that all zeros are eventually simple and negative as the number of rows
increases, and that successive polynomials then strictly interlace.
A discrete sine comparison identifies two dominant diagonal points of a
torus integral. Uniform saddle estimates and the absolute indexing of
zeros in two entire endpoint limits account for every zero.
\end{abstract}

\begin{theorem}\label{thm:main}
Fix integers \(d\ge2\) and \(1\le r<d\). Let \(a_{n,k}\) count labeled
\(n\times d\) binary arrays with each column sum equal to \(k\) and each
row sum at most \(r\). Put
\[
P_n(q)=\sum_k a_{n,k}q^k,\qquad D_n=\lfloor nr/d\rfloor.
\]
There is \(N(d,r)\) such that for \(n\ge N(d,r)\) the \(D_n\) zeros of
\(P_n\) are simple and negative. The pair \(P_n,P_{n+1}\) strictly
interlaces, with the rightmost zero belonging to \(P_{n+1}\), whenever
there are zeros to compare.
\end{theorem}

The balanced-array conjecture asks for real-rootedness for every
\(n,d,r\) \cite{balanced}. Its discussion reports a preliminary local
eventual argument at \(d=4,r=2\) based on endpoint limits and saddle
estimates. Here the threshold is existential for each fixed pair
\((d,r)\). The result gives a large-row theorem throughout the proper
row caps. It leaves the general finite gap to the all-row conjecture.

For \(n\ge1\), every integer \(k\) from \(0\) to \(D_n\) is feasible.
Take \(dk\) consecutive positions in a cyclic list of the \(n\) row
labels and use successive blocks of \(k\) positions to mark the ones in
successive columns. Each block uses distinct labels, and every row
receives at most \(\lceil dk/n\rceil\le r\) ones. This proves the
degree assertion together with the necessary inequality \(dk\le nr\).

When \(d=2,r=1\), direct enumeration gives
\(a_{n,k}=n!/((n-2k)!(k!)^2)\).
The generating function in \(n\) is
\[
\sum_{n\ge0}P_n(q)x^n=\big(1-2x+(1-4q)x^2\big)^{-1/2},
\]
by the constant term of \((1+\sqrt q(z+z^{-1}))^n\).
Thus for \(q<0\),
\[
P_n(q)=(1-4q)^{n/2}L_n\big((1-4q)^{-1/2}\big),
\]
where \(L_n\) is the Legendre polynomial \cite{dlmf}.
Its positive zeros are simple and correspond to negative zeros
under the increasing map \(y\mapsto(1-y^{-2})/4\).
Parity gives exactly \(\lfloor n/2\rfloor\) positive zeros of \(L_n\).
The interlacing of successive Legendre polynomials, restricted to
their positive zeros, gives the desired directed interlacing.
In the rest of the proof assume \(d\ge3\).

\section{Endpoint limits}

We first establish the limiting behavior needed at the ends.
Suppress the subscript in \(D=D_n\), and put
\[
s=nr-dD,\quad
\gamma=\frac{\binom d{r-1}}{\binom dr}=\frac{r}{d-r+1},\quad
E_d(x)=\sum_{j\ge0}\frac{x^j}{(j!)^d},\quad
H_s(x)=\sum_{j\ge0}\frac{x^j}{(dj+s)!}.
\]

\begin{lemma}\label{lem:limits}
The following limits hold locally uniformly in \(x\in\mathbb C\),
with \eqref{eq:reverse} taken along the residue \(s\).
\begin{align}
P_n(x/n^d)&\longrightarrow E_d(x),\label{eq:forward}\\
n\big(P_{n+1}(x/n^d)-P_n(x/n^d)\big)&\longrightarrow d x E_d'(x),
\label{eq:difference}\\
\sum_{j=0}^{D}\frac{a_{n,D-j}}{a_{n,D}}
      \left(\frac{x}{(\gamma n)^d}\right)^j&\longrightarrow s!H_s(x).
\label{eq:reverse}
\end{align}
\end{lemma}

\begin{proof}
The unrestricted count \(b_{n,k}=\binom nk^d\) is obtained by choosing
the column supports independently. A union bound over rows and over
sets of \(r+1\) columns gives, for \(0\le k\le n\),
\[
0\le b_{n,k}-a_{n,k}
 \le n\binom d{r+1}(k/n)^{r+1} b_{n,k}.
\]
For fixed \(k\), \(a_{n,k}/n^{dk}\to1/(k!)^d\), and this expression is
bounded by \(1/(k!)^d\) for all \(k\). This proves \eqref{eq:forward}.
At fixed \(k\), the unrestricted differences satisfy
\(n(b_{n+1,k}-b_{n,k})/n^{dk}\to dk/(k!)^d\) by the binomial ratio.
Pascal's identity bounds the scaled differences by
\(e^{2d}dk/(k!)^d\) on \(k\le n+1\), since
\(b_{n+1,k}-b_{n,k}\le dk(n+1)^{dk-1}/(k!)^d\).
For \(r\ge2\), the two capped-to-unrestricted errors after the same
scaling vanish at fixed \(k\), and the union bound dominates them
by a constant times \(k^{r+1}/(k!)^d\). This proves
\eqref{eq:difference} in that case.
For \(r=1\) the column supports must be disjoint, so
\(a_{n,k}=(n)_{dk}/(k!)^d\) with a falling factorial in the numerator.
The constant coefficient is fixed, and for \(k\ge1\) the identity
\((n+1)_{dk}-(n)_{dk}=dk(n)_{dk-1}\) proves the same limit with
majorant \(dk/(k!)^d\).

We give the lattice estimate used for \eqref{eq:reverse}.
Write \(h=d-1\), and let \(Y\in\mathbb Z^h\) be the indicator vector
of the first \(h\) positions of a uniformly random \(r\)-subset of
\([d]\). Write \(\iota=(1,\ldots,1)^\top\) for the \(h\)-dimensional ones vector.
The mean is \(\mu=(r/d)\iota\), and the covariance is
\[
\Sigma=\frac{r(d-r)}{d(d-1)}
       \left(I_h-\frac1d\iota\iota^\top\right)>0.
\]
Indeed the diagonal variances are \(r(d-r)/d^2\), and the off-diagonal
covariances are \(-r(d-r)/(d^2(d-1))\). The support differences
include each standard basis vector, by fixing \(r-1\) positions
outside \(\{i,d\}\) and adjoining either \(i\) or \(d\).
Thus the characteristic function \(\varphi(t)=\mathbb E e^{it\cdot Y}\)
has modulus one on \([-\pi,\pi]^h\) only at \(0\).
For \(S_M\) a sum of \(M\) independent copies of \(Y\), Fourier
inversion and the centered expansion
\[
e^{-it\cdot\mu}\varphi(t)=1-\frac12 t^\top\Sigma t+O(\|t\|^3)
\]
give
\begin{align*}
\sup_v\Pr(S_M=v)&=O_{d,r}((M+1)^{-h/2}),\\
\Pr(S_M=v)&=M^{-h/2}\big(\kappa+o(1)\big)
   \quad(\|v-M\mu\|\le B),
\end{align*}
where \(v\) ranges over \(\mathbb Z^h\), the last estimate is
uniform for bounded \(B\), and
\(\kappa=(2\pi)^{-h/2}(\det\Sigma)^{-1/2}\).
For the upper bound use \(|\varphi(t)|\le e^{-c\|t\|^2}\)
near zero and a strict compact modulus gap on the rest of the
cube. For the asymptotic, rescale the integral by \(t=w/\sqrt M\)
in the small neighborhood. The phase from \(v-M\mu\) tends
uniformly to 1 on compact sets of \(w\), the centered power
converges to \(e^{-w^\top\Sigma w/2}\), and Gaussian domination
controls the tails. This also fixes the lattice factor.

Set \(A=\binom dr\) and \(b_i=\binom d{r-i}/A\) for \(1\le i\le r\).
A row of size \(r-i\) has deficiency \(i\). At column sum \(D-j\),
the total deficiency is \(e=s+dj\). Without balance the count is
\[
U_{n,e}=A^n[z^e]\left(1+\sum_{i=1}^{r}b_i z^i\right)^n.
\]
For fixed \(e\) the leading term comes from \(e\) rows of deficiency
one, hence \(U_{n,e}\sim A^n n^e\gamma^e/e!\).
Conditional on the ordered row sizes, choose rows independently
and uniformly among subsets of those sizes. There are \(M\)
regular rows of size \(r\), with \(M=n-O(e)\). Given all deficient
rows, the targets of the first \(h\) projected regular column
counts differ from \(M\mu\) by \(O(e)\), since
\((D-j)-nr/d=-e/d\) and at most \(e\) rows are deficient.
Once those \(h\) targets are reached, the last column sum follows
from the fixed total. At fixed \(e\) the Fourier asymptotic thus
gives
\[
a_{n,D-j}\sim \kappa A^n n^{e-h/2}\gamma^e/e!.
\]
It follows that the coefficient of \(x^j\) in the reverse scaling
tends to \(s!/(s+dj)!\).

To dominate the coefficients uniformly in \(j\), put
\(W_e=[z^e]\exp(\sum_{i=1}^r b_i z^i)\).
Choosing the deficient rows first gives \(U_{n,e}\le A^n n^e W_e\):
there are at most \(e\) deficient rows, and
\(\binom n v\le n^v/v!\) for their number \(v\).
If \(e\le n/2\), the Fourier upper bound supplies a factor
\(O_{d,r}(n^{-h/2})\) on the conditional balance probability, since
\(M\ge n/2\). If \(e>n/2\), use probability at most one, and
\(n^{h/2}\le(2e)^{h/2}\).
A uniform lower bound for the top coefficient follows from its
asymptotic with \(e=s\), since only finitely many residues occur.
Altogether the coefficient at \(j\) in \eqref{eq:reverse} is at
most
\[
C_{d,r}(1+e)^{h/2}\gamma^{-dj}W_e,\qquad e=s+dj.
\]
For each \(R>0\), \(W_e\le R^{-e}\exp(\sum_i b_i R^i)\).
Choosing \(R\) sufficiently large gives a summable majorant on
any specified disk in \(x\). This proves \eqref{eq:reverse}.
\end{proof}

\begin{lemma}\label{lem:simple}
The functions \(E_d,H_0,\ldots,H_{d-1}\) have infinitely many
simple negative zeros. If \(\eta_{s,i}\) is the \(i\)-th zero
of \(H_s\) from the right, then
\[
\eta_{0,i}>\eta_{1,i}>\cdots>\eta_{d-1,i}>\eta_{0,i+1}.
\]
\end{lemma}

\begin{proof}
The negative zero location is a special case of the hyper-Bessel
zero theorem \cite[Theorem 2]{hyperbessel}. In generalized
hypergeometric notation, \(E_d\) is \({}_0F_{d-1}\) with all
denominator parameters one. The function \(s!H_s\) has
parameters \((s+1)/d,\ldots,(s+d)/d\) with the parameter one
removed, at argument \(x/d^d\). Their series give order at
most \(1/d<1\). For instance, the absolute series for \(E_d\)
is bounded by \((\sum_{j\ge0}|x|^{j/d}/j!)^d\), and that for
\(H_s\) by \(e^{|x|^{1/d}}\) for \(|x|\ge1\).
Hadamard factorization gives infinitely many zeros and a
genus-zero product because these functions are nonpolynomial.

Put \(\vartheta=x\partial_x\). The coefficient identities are
\[
\vartheta^d E_d=xE_d,\qquad
\prod_{\nu=0}^{d-1}(d\vartheta+s-\nu)H_s=xH_s.
\]
Each intermediate function in either operator product has
only real nonpositive zeros. Indeed canonical product
truncations approximate the initial functions locally uniformly
by polynomials with such zeros. The operators \(\vartheta\)
and \(\vartheta+\alpha\) for \(\alpha>0\) preserve this property.
The former follows from Rolle's theorem. For the latter,
on simple negative roots the values \(xf'(x)\) alternate, and
the sign at zero gives the final sign change, accounting for
the degree. Approximation handles multiple roots and zeros at zero.
In the product for \(H_s\), apply \(d\vartheta\) first, obtaining
a factor \(x\). Each remaining factor on the quotient is
\(d\vartheta+d+s-\nu\) with positive constant.
The operators pass to locally uniform limits. Their intermediate
functions are nonpolynomial of order below one by their series.

If \(f\) has a zero of multiplicity \(h_0\ge1\) at a negative
point, \(d\vartheta+\alpha\) lowers it to \(h_0-1\).
If \(f\) is nonzero there, use its genus-zero product to obtain
\[
\frac{\dd}{\dd x}\left(\frac{f'}{f}+\frac{\alpha}{dx}\right)
=-\sum_{\rho:f(\rho)=0}\frac1{(x-\rho)^2}-\frac{\alpha}{dx^2}.
\]
This is strictly negative for \(\alpha\ge0\), and also if \(f\)
already vanishes at zero and \(\alpha>-d\), using that zero's
contribution. Thus any newly introduced negative zero is simple
for the factors above. This also applies to \(\vartheta\) by
scaling. An initial multiplicity at least two would strictly
decrease across the product, contradicting its preservation
by the right-hand side \(xf\). This proves simplicity.

For \(s\ge1\), \((d\vartheta+s)H_s=H_{s-1}\).
The logarithmic derivative of \(H_s\) plus \(s/(dx)\)
goes strictly from positive infinity to negative infinity
in each gap and in the rightmost interval ending at zero.
The unmodified logarithmic derivative of \(H_0\) is positive
on its rightmost interval, and \(dH_0'=H_{d-1}\) gives the
closing comparison. These facts yield the displayed order.
\end{proof}

\section{The two dominant points of the torus integral}

Put \(\phi=\pi/d\), \(\zeta=e^{i\phi}\), \(a_j=\binom dj\), and
\[
R(t,y)=\sum_{j=0}^r t^j e_j(y),\qquad
F(z)=\sum_{j=0}^r a_j z^j,\qquad S(t)=F(t\zeta).
\]
Here \(e_j\) is the \(j\)-th elementary symmetric polynomial.
Integrate on the torus of unit coordinates \(y_i\) with product
\(-1\), using independent uniform angles for the first \(d-1\)
coordinates and \(y_d=-1/\prod_{i<d}y_i\).
Then for \(t>0\),
\begin{equation}\label{eq:torus}
P_n(-t^d)=\int R(t,y)^n\,\dd y,
\end{equation}
where the measure has mass one. Indeed this constant term
selects equal exponents \(k\) in all coordinates and gives
factor \((-t^d)^k\).

\begin{lemma}\label{lem:dominance}
For every \(t>0\) the modulus of \(R\) has exactly two
maximizing points on this torus, the all-\(\zeta\) point and
its conjugate. Both are nondegenerate maxima of \(\log|R|\).
\end{lemma}

\begin{proof}
First, for integers \(1\le j<d\) and \(1\le k\le d\),
\begin{equation}\label{eq:sines}
\left|\frac{\sin(kj\phi)}{\sin(j\phi)}\right|
 \le \frac{\sin(k\phi)}{\sin\phi}.
\end{equation}
Both sides are zero if \(k=d\). Otherwise use the symmetries
\(j\mapsto d-j\), \(k\mapsto d-k\) to take \(j,k\le d/2\).
The case \(k=1\) is immediate. For \(k\ge2\) and
\(x=j\phi\le\pi/k\), the nonnegative ratio
\(\sin(kx)/\sin x\) decreases on \((0,\pi/k)\), since its
logarithmic derivative is \(k\cot(kx)-\cot x<0\) by the
strict decrease of \(u\cot u\) on \((0,\pi)\). For
\(\pi/k\le x\le\pi/2\), the ratio in absolute value is at
most \(1/\sin(\pi/k)\le k/2\), while
\(\sin(k\phi)/\sin\phi\ge2k/\pi>k/2\).
This proves \eqref{eq:sines}.

Consider the diagonal values \(F(te^{i\alpha})\) with
\(\alpha\) an odd multiple of \(\phi\). In the coefficient
of \(t^s\) in \(|F(te^{i\alpha})|^2\), the frequencies are
\(2i-s\) with weights \(\binom di\binom d{s-i}\) over
\(\max(0,s-r)\le i\le\min(s,r)\).
The weights are symmetric and increase to the center, by
the consecutive binomial ratios. Decompose them into
nonnegative layers on nested centered intervals. The
frequencies on a layer have form \(-m,-m+2,\ldots,m\),
with \(m\le r\), so their exponential sum is
\(\sin((m+1)\alpha)/\sin\alpha\).
At \(\alpha=j\phi\), \(1\le j<d\), \eqref{eq:sines}
bounds every coefficient by its value at \(\phi\).
For \(j>1\) the coefficient of \(t\), namely
\(2d\cos(j\phi)\), is strictly smaller. Conjugation
handles the negative angles.

There remains \(\alpha=\pi\) when \(d\) is odd.
Let \(u_i=a_i t^i\) for \(0\le i\le r\) and choose a peak
\(p\). Consecutive ratios \(u_{i+1}/u_i\) strictly decrease
when two such ratios exist. Positive superlevel sets
are intervals containing \(p\). On any such interval the
alternating sum at \(\pi\) has absolute value at most one,
so \(|F(-t)|\le u_p\) by summing the threshold layers.
For the evaluation at \(\phi\), project each interval's
geometric sum onto the direction \(e^{ip\phi}\).
For length \(L\le d\) this projection is minimized when
the peak is an endpoint of the interval. Its minimum is
\[
\frac{\sin(L\phi/2)\cos((L-1)\phi/2)}{\sin(\phi/2)}
 =\frac12\left(1+\frac{\sin((2L-1)\phi/2)}{\sin(\phi/2)}\right)
 \ge1.
\]
The bound is strict if \(1<L<d\), or if \(L=d\) and \(p\)
is interior. If \(r+1<d\), the full support gives strictness.
If \(r+1=d\) and an interior peak exists, choose it.
Otherwise the unique peak is an endpoint, and the strict
ratio inequalities force a layer of length two with
positive weight. Since odd \(d\ge3\), this layer also gives
strictness. Thus \(|S(t)|>u_p\ge|F(-t)|\).
The only diagonal maximizers are the nearest conjugate pair.

To lift the comparison to the torus, take any connected
component of the compact maximizing set, and maximize
\(|\sum_i y_i|\) on it. If \(y_i\ne y_j\) at this point,
hold their product and all other coordinates fixed. Their
sum ranges continuously on the full segment from
\(-2\sqrt{y_i y_j}\) to \(2\sqrt{y_i y_j}\), with current
sum in its interior. Symmetry and multiaffinity make \(R\)
affine in the pair sum. Unless the slope is zero, strict
convexity of the squared modulus excludes an interior
maximum. If the slope vanishes, the full pair variation
stays in the same maximizing component, while
\(|\sum_i y_i|^2\) is strictly convex along the segment.
This contradicts the auxiliary maximum. Every maximizing
component therefore contains a diagonal point, which must
be one of the two preferred ones.

At the all-\(\zeta\) point put \(z=t\zeta\) and
\(B(z)=\sum_{\ell=0}^{r-1}\binom{d-2}{\ell}z^\ell\).
Its terms lie in a sector of width \((r-1)\phi<\pi\), so
\(B(z)\ne0\). Perturbing a pair of angles by \(p,-p\)
changes \(S\) to \(S+2zB(z)(\cos p-1)\).
Maximality compared to \(p=\pi\) gives
\[
\Re(zB(z)\overline S)\ge 2|zB(z)|^2>0.
\]
In angular displacements \(p_i\) with \(\sum_i p_i=0\),
write \(Q=\sum_i p_i^2\). Expanding each elementary
symmetric polynomial gives
\begin{equation}\label{eq:quadratic}
\log(R/S)=-\lambda(t)Q/2+O(\|p\|^3),
 \qquad \lambda(t)=\frac{zB(z)}{S(t)}.
\end{equation}
Indeed \(\sum_{|I|=j}(\sum_{i\in I}p_i)^2
=\binom{d-2}{j-1}Q\). The branch of the logarithm
vanishes at the point. The displayed real-part inequality
gives \(\Re\lambda>0\), so both diagonal maxima are
nondegenerate and isolated. Each maximizing component
contains one of these and hence is a singleton.
\end{proof}

\section{Uniform overlap estimates}

Put \(\nu=(d-1)/2\),
\(\theta(t)=\Arg S(t)\),
\(\Psi_n(t)=n\theta(t)-\nu\Arg\lambda(t)\), and
\[
C_n(t)=\frac{\sqrt d\,(2\pi n)^\nu|\lambda(t)|^\nu}{2|S(t)|^n}.
\]
Use the arguments \(\theta\in(0,r\phi)\),
\(\Arg\lambda\in(-\pi/2,\pi/2)\).
All terms of \(S\) lie in a sector of width \(r\phi<\pi\),
and \(\Im S>0\).

\begin{lemma}\label{lem:overlap}
For a common fixed sufficiently large \(K\), and all sufficiently
large \(n\), the following two estimates hold on
\(K/n\le t\le\gamma n/K\):
\begin{align}
 |C_n P_n(-t^d)-\cos\Psi_n|&<\tfrac12,\label{eq:sign}\\
 \left|\frac{C_n}{\Im S}
 \big(P_{n+1}(-t^d)-(\Re S)P_n(-t^d)\big)+\sin\Psi_n\right|
 &<\tfrac12.\label{eq:sw}
\end{align}
The phase is strictly increasing there.
\end{lemma}

\begin{proof}
The two integral amplitudes in \eqref{eq:torus} are
\(1\) and \((R-\Re S)/\Im S\). The latter has values \(i,-i\)
at the conjugate saddles. In the \(d-1\) free coordinates,
the matrix of \(Q\) is \(I+\iota\iota^\top\), of
determinant \(d\). At the first saddle the leading term
of the integral is
\[
\frac{S^n}{\sqrt d\,(2\pi n)^\nu\lambda^\nu}
\]
times the amplitude there. The power of \(\lambda\) uses
the indicated argument, from analyticity of the Gaussian
integral in the right half-plane. The other leading term
is conjugate. On any fixed compact \(t\)-interval the
quadratic form has a uniform strictly negative real part
in \eqref{eq:quadratic}, and there is a uniform modulus
gap off fixed small neighborhoods. Interpolating between
the quadratic exponent and the full exponent bounds
the difference of their \(n\)-th exponentials in a
normalized neighborhood by
\(O(n\|p\|^3 e^{-cn\|p\|^2})\). Integrating costs
\(O(n^{-\nu-1/2})\), also the size of the error from
variation of a uniformly smooth amplitude. Off the
neighborhoods the contribution is exponentially smaller.

As \(t\downarrow0\), \(\lambda/t\to\zeta\) with
\(\Re\zeta>0\) because \(d\ge3\). The cubic remainder
in \eqref{eq:quadratic} is \(O(t\|p\|^3)\), uniformly
since the normalized logarithm vanishes identically
at \(t=0\). The torus function \(\Re e_1(y)\)
has maximum \(d\cos\phi\) only at the same conjugate pair.
Indeed the component argument of Lemma~\ref{lem:dominance}
applies to this real affine function of each pair sum.
If its slope along the segment is zero the auxiliary
strict convexity applies as before. The two maximizing
diagonal points are isolated by their negative real
quadratic coefficient \(-(\cos\phi)Q/2\).
The expansions \(|R|^2=1+2t\Re e_1+O(t^2)\) and
\(|S|^2=1+2td\cos\phi+O(t^2)\) thus give
\(|R|/|S|\le e^{-ct}\) off those neighborhoods.
The weighted amplitude and its angular derivatives are
uniformly bounded since both its numerator and denominator
are of order at most \(t\), with the denominator of exact
order \(t\). Gaussian scaling with parameter \(nt\) and
the complement bound give normalized errors
\[
O\big((nt)^{-1/2}+(nt)^\nu e^{-cnt}\big).
\]

At infinity use \(v=1/t\) and divide \(R,S\) by \(t^r\).
At \(v=0\), the modulus of \(e_r\) reaches \(a_r\) only
at the \(d\) diagonal points. Equality requires all
\(r\)-subset products to agree, and comparison of subsets
differing only by exchanging \(i,j\) gives \(y_i=y_j\).
Such subsets exist since \(1\le r<d\).
Consequently there is a uniform strict modulus gap off
fixed neighborhoods of all \(d\) points. At each point
the normalized quadratic coefficient tends to
\(\binom{d-2}{r-1}/a_r>0\). At a diagonal angle \(\alpha\)
the rescaled modulus is
\(a_r+a_{r-1}v\cos\alpha+O(v^2)\).
Thus the additional points have modulus ratio at most
\(e^{-cv}\) to \(|S/t^r|\) for small \(v>0\), and their
local integrals have size bounded by their modulus to
the \(n\)-th power times \(O(n^{-\nu})\).
The weighted amplitude stays uniformly smooth after
division of numerator and denominator by \(t^r\),
since \(t^{-r}\Im S\) has a positive limit.
This yields normalized errors
\[
O(n^{-1/2}+e^{-cnv}+n^\nu e^{-cn}).
\]
Combining the three ranges and selecting \(K\) large
enough for both amplitudes gives the sign errors
less than \(1/2\) on the overlap. The conjugate leading
terms give \(\cos\Psi_n\) and
\(\Re(i e^{i\Psi_n})=-\sin\Psi_n\).

Finally
\[
|S|^2\theta'(t)=
\sum_{0\le\ell<j\le r}(j-\ell)a_j a_\ell t^{j+\ell-1}
                \sin((j-\ell)\phi)>0.
\]
Near zero the ratio \((\Arg\lambda)'/\theta'\) is bounded
because \(\lambda/t\) extends smoothly and nonvanishingly,
and \(\theta'\) tends to \(d\sin\phi>0\).
Near infinity use \(v=1/t\), where \(\lambda\) has a
smooth nonzero extension, and
\(\theta(1/v)=r\phi-\gamma\sin\phi\,v+O(v^2)\).
Both derivatives in \(t\) are of size at most \(O(t^{-2})\),
and \(\theta'\) has a positive leading coefficient of that
order. Compactness handles the intermediate range.
Therefore \(\Psi_n'\) is positive for all sufficiently
large \(n\).
\end{proof}

\section{Absolute indexing of the endpoint zeros}

Write \(b=\sin\phi\), \(c=db\), and let \(T_i\) and \(T_{s,i}\)
be the positive \(d\)-th roots of the magnitudes of the zeros
of \(E_d\) and \(H_s\), respectively, in increasing order.
The index shifts in their asymptotics will fix the integer
phase cutoffs.

\begin{lemma}\label{lem:index}
For fixed \(d\ge3\) and \(0\le s<d\),
\begin{align}
T_i&=\frac{\pi}{c}(i+\nu/d-\tfrac12)+O(i^{-1/2}),
\label{eq:index-small}\\
T_{s,i}&=\frac{\pi}{b}(i+s/d-\tfrac12)+O(e^{-c_1i})
\label{eq:index-large}
\end{align}
for some \(c_1>0\).
\end{lemma}

\begin{proof}
We first record an index calculation for an entire function
\(f\) of order below one with \(f(0)>0\) and only simple
negative roots \(-V_i^d\) in increasing magnitude.
Suppose \(V_i=a(i+\delta)+O(i^{-\epsilon})\) for
\(a,\epsilon>0\). By the canonical product,
\(\log f(T^d)=\log f(0)+\sum_i\log(1+(T/V_i)^d)\).
When \(V_i\) is replaced by \(a(i+\delta)\) for \(i\ge I\)
with \(I+\delta>0\), the error is uniformly bounded in \(T\),
since the derivative of each term with respect to \(V_i\)
has absolute value at most \(d/V_i\).

Put \(\beta=I+\delta\), \(U=T/a\), and let \(\alpha\)
run over the roots of \(\alpha^d=(-1)^{d+1}\).
The Euler gamma limit and \(\sum_\alpha\alpha=0\) give
\[
\prod_{j\ge0}\left(1+\frac{U^d}{(\beta+j)^d}\right)
=\prod_{\alpha^d=(-1)^{d+1}}
   \frac{\Gamma(\beta)}{\Gamma(\beta+\alpha U)}.
\]
Indeed each factor on the left splits into
\(\prod_\alpha(1+\alpha U/(\beta+j))\), while the
normalizing powers in the gamma limits cancel.
No \(\alpha\) is negative real, so Stirling's formula
applies in sectors avoiding the negative axis \cite{dlmf}.
Its real logarithm gives
\[
-U\Re\sum_\alpha\alpha\log\alpha
 -d(\beta-\tfrac12)\log U+O(1).
\]
For the principal arguments \(\vartheta_j=(2j-d+1)\phi\),
\(0\le j<d\), the geometric sum
\(\sum_j e^{i\vartheta_j x}=\sin(\pi x)/\sin(\phi x)\)
has derivative \(-\pi/\sin\phi\) at \(x=1\).
Thus \(\Re\sum_\alpha\alpha\log\alpha=-\pi/\sin\phi\).
The initial \(I-1\) factors contribute \(d(I-1)\log T+O(1)\),
giving
\begin{equation}\label{eq:offset-product}
\log f(T^d)=\frac{\pi T}{a\sin\phi}
            -d(\delta+\tfrac12)\log T+O(1).
\end{equation}

The torus average of \(\exp(T e_1(y))\) with product
\(-1\) is \(E_d(-T^d)\), by expanding the exponentials and
taking the constant term. Its real exponent has the same
two strict maximizing points as in the small-\(t\) estimate.
The expansion at the first point is
\(e_1=d\zeta-\zeta Q/2+O(\|p\|^3)\).
The Gaussian estimate gives
\[
G(T):=\frac{\sqrt d\,(2\pi T)^\nu}{2}
       e^{-dT\cos\phi}E_d(-T^d)
 =\cos(cT-\nu\phi)+O(T^{-1/2}).
\]
Differentiating the integral introduces the amplitude
\(e_1\), with value \(d\zeta\) at the first point.
The real term \(d\cos\phi\) cancels against the derivative
of the exponential normalization, and the term \(\nu/T\)
is within the error. Thus
\(G'(T)=-c\sin(cT-\nu\phi)+O(T^{-1/2})\).
These estimates place exactly one zero in each sufficiently
large fixed-width band around a zero of the cosine, and no
zeros outside them. On each band the derivative has fixed
sign and is bounded away from zero, and the root differs
from its band center by \(O(T^{-1/2})\).
This gives \(T_i=(\pi/c)(i+\delta)+O(i^{-1/2})\) for some
fixed index shift \(\delta\).

For the positive argument use the same exponential integral
on the torus with product \(+1\).
The real part of \(e_1\) is maximized only when all coordinates
are one, with quadratic coefficient one.
Therefore \(\log E_d(T^d)=dT-\nu\log T+O(1)\).
Comparing the logarithmic term in \eqref{eq:offset-product}
gives \(\delta=\nu/d-\tfrac12\), proving
\eqref{eq:index-small}.

For \(H_s\) the exponential root-of-unity filter gives
\[
dT^s H_s(-T^d)
 =\sum_{\eta^d=-1}\eta^{-s}e^{\eta T}
 =2e^{T\cos\phi}\cos(bT-s\phi)
   +O(e^{T\cos(3\phi)}).
\]
The two leading conjugate exponentials have real part
\(T\cos\phi\), and the normalized error and its derivative
are exponentially small. The band argument therefore gives
\(T_{s,i}=(\pi/b)(i+\delta_s)+O(e^{-c_1i})\).
Filtering with roots of \(+1\) gives the single leading
exponential and \(\log H_s(T^d)=T-s\log T+O(1)\).
Equation \eqref{eq:offset-product} yields
\(\delta_s=s/d-\tfrac12\), proving \eqref{eq:index-large}.
\end{proof}

\section{Matching the phases and interlacing}

At zero, \(\theta(t)=ct+O(t^2)\) and
\(\Arg\lambda(t)\to\phi\).
At infinity, \(\theta(t)=r\phi-\gamma b/t+O(t^{-2})\)
and \(\Arg\lambda(t)\to0\).
Hence for \(T\) in a compact positive interval
\begin{equation}\label{eq:phase-limit}
\Psi_n(T/n)\longrightarrow cT-\nu\phi,\qquad
\Psi_n(\gamma n/T)-D\pi\longrightarrow s\phi-bT,
\end{equation}
with the second limit along the residue \(s\), since
\(nr\phi-D\pi=s\phi\).

Choose \(K\) as in Lemma~\ref{lem:overlap}, for both estimates.
Choose fixed large positive integers \(L,M\) with
\[
U=\frac{\pi L+\nu\phi}{c}>K,\qquad
V_s=\frac{\pi M+s\phi}{b}>K.
\]
The index formulas show that \(U\) falls strictly between
\(T_L,T_{L+1}\) and that each \(V_s\) falls strictly between
\(T_{s,M},T_{s,M+1}\) once \(L,M\) are sufficiently large.
The two gaps each have a positive asymptotic half-spacing
on either side.

For all sufficiently large \(n\) there are unique nodes
\(t_{\ell,n}<t_{h,n}\) in the overlap satisfying
\(\Psi_n(t_{\ell,n})=L\pi\) and
\(\Psi_n(t_{h,n})=(D-M)\pi\).
Indeed bracket each endpoint limit by two fixed \(T\)
values greater than \(K\), use \eqref{eq:phase-limit}
with opposite signs on the brackets, and apply strict
monotonicity. Shrinking the brackets yields
\(nt_{\ell,n}\to U\) and
\(\gamma n/t_{h,n}\to V_s\) in the indicated residues.
The order of the nodes follows from \(D-M>L\) for large \(n\).

In the forward scaling of Lemma~\ref{lem:limits}, choose
two circles about zero with radii bracketing \(U^d\) in
the same gap between consecutive limit root magnitudes.
Rouch\'e's theorem on both circles gives exactly \(L\)
roots with \(|q|<t_{\ell,n}^d\), and on disjoint small
conjugation-invariant disks centered at the negative limit
zeros it gives one root in each disk. Thus these roots are
simple and negative. In the reversed scaling use radii
bracketing \(V_s^d\), where a nonzero root \(q\) corresponds
to \(x=(\gamma n)^d/q\). This gives exactly \(M\) simple
negative roots with \(|q|>t_{h,n}^d\).
Finiteness of the residues makes the assertions simultaneous.

For every integer \(j=L,\ldots,D-M\), the intermediate
node of phase \(j\pi\) has sign \((-1)^j\) by \eqref{eq:sign}.
Between successive nodes there are therefore \(D-M-L\)
distinct negative roots. They are disjoint from the endpoint
regions, and the total count is \(D\). This proves the
real-rootedness and simplicity assertion of
Theorem~\ref{thm:main}.

We next compare successive polynomials. At a central
root \(q=-t^d\) with phase in \((j\pi,(j+1)\pi)\),
\eqref{eq:sign} gives \(|\sin\Psi_n|>\sqrt3/2\).
The weighted estimate \eqref{eq:sw} shows that
\(P_{n+1}(q)\) has sign opposite to \((-1)^j\), the sign
of this sine. There is exactly one simple root in the
node interval, where \(P_n(-t^d)\) crosses from sign
\((-1)^j\) to sign \((-1)^{j+1}\) as \(t\) increases.
Since \(q=-t^d\) decreases with \(t\), \(P_n'(q)\)
has sign \((-1)^j\). Thus \(P_{n+1}(q)/P_n'(q)<0\).

We show the same for the fixed number of endpoint roots.
Let \(R_{n,i}\) denote the \(i\)-th negative root from the
right, and let \(\xi_i\) be the corresponding zero of \(E_d\).
By \eqref{eq:forward}, \(X_n=n^dR_{n,i}\) and
\(Y_n=n^dR_{n+1,i}\) tend to \(\xi_i\).
The difference limit \eqref{eq:difference} evaluated at
\(X_n\) gives
\(nP_{n+1}(X_n/n^d)\to d\xi_i E_d'(\xi_i)\).
The derivative of \(x\mapsto P_{n+1}(x/n^d)\) tends
locally uniformly to \(E_d'\), by \eqref{eq:forward}
and the argument scaling \((1+1/n)^d x\).
The mean value theorem applied between \(X_n,Y_n\)
therefore gives \(n(X_n-Y_n)\to d\xi_i<0\).
Distinct limits also give the next-index gap, so
\[
R_{n+1,i}>R_{n,i}>R_{n+1,i+1}
\]
eventually at each fixed index.

At the left, let \(L_{n,i}\) denote the \(i\)-th root
from the left and let \(s'=(s+r)\bmod d\) be the next
residue. Lemma~\ref{lem:simple} implies
\(\eta_{u,i}>\eta_{v,i}>\eta_{u,i+1}\) for every \(u<v\),
using the closing comparison with \(\eta_{0,i+1}\).
Reciprocate the limit roots in \eqref{eq:reverse},
noting that the scale ratio \((1+1/n)^d\) tends to one.
If the residue wraps, then \(s'<s\), the degree rises by
one, and \(L_{n+1,i}<L_{n,i}<L_{n+1,i+1}\) eventually.
Otherwise \(s'>s\), the degree is unchanged, and
\(L_{n,i}<L_{n+1,i}<L_{n,i+1}\) eventually.

At the \(i\)-th root from the right, \(P_{n+1}\) has
exactly \(i\) roots to its right, while \(P_n'(q)\) has
sign \((-1)^{i-1}\).
At the \(i\)-th root from the left, \(P_{n+1}\) has
exactly \(i\) roots to its left if the degree increases,
and \(i-1\) otherwise. In either case the signs from
the positive leading coefficients give the same negative
quotient as in the center. All comparisons here are
needed only for finitely many fixed indices, together
with one more from each end.

Consequently \(P_{n+1}\) has alternating signs at
successive roots of \(P_n\) and a root between the
rightmost root and zero, since its constant coefficient
is positive. If the degree increases, the sign at negative
infinity supplies a further root to the left. These signs
account for its entire degree and prove the directed
strict interlacing.

\begin{thebibliography}{9}
\bibitem{balanced}
P. Alexandersson, \emph{Open problems on real-rooted polynomials},
SymCat. See the balanced-array conjecture and
\href{https://github.com/PerAlexandersson/RealRooted/issues/2005}
{RealRooted issue \#2005}.

\bibitem{hyperbessel}
\'A. Baricz and S. Singh, Zeros of some special entire functions,
\href{https://arxiv.org/abs/1702.00626}{arXiv:1702.00626},
Theorem 2.

\bibitem{dlmf}
\emph{NIST Digital Library of Mathematical Functions}.
See the \href{https://dlmf.nist.gov/18.12.E11}{Legendre generating
function}, \href{https://dlmf.nist.gov/18.2.vi}{orthogonal-polynomial
zero interlacing}, \href{https://dlmf.nist.gov/5.8.E1}{Euler gamma limit},
and \href{https://dlmf.nist.gov/5.11.E1}{Stirling expansion}.
\end{thebibliography}

\end{document}
